\documentclass[aps,prd,twocolumn,superscriptaddress,nofootinbib]{revtex4-2}

\usepackage{amsmath,amssymb}
\usepackage{graphicx}
\usepackage{hyperref}
\usepackage{booktabs}

\begin{document}

\title{The Sorkin-Johnston vacuum on causal sets and causal dynamical triangulations:\\a cross-approach comparison}

\author{Matt Loftus}
\email{matthew.a.loftus@gmail.com}
\affiliation{Independent researcher}

\date{\today}

\begin{abstract}
We compute the Sorkin-Johnston vacuum for a free massless scalar field on two independent discretizations of 2D spacetime---causal sets (random 2-orders) and causal dynamical triangulations (CDT)---using identical numerical methods. CDT reproduces the expected central charge of $c = 1$ for a free scalar ($c_{\mathrm{eff}} = 1.36 \to 0.98 \to 1.15$ at $N = 38$--$127$), while causal sets yield $c_{\mathrm{eff}} \approx 3$ at small $N$ and a \emph{diverging} effective central charge at large $N$ ($c_{\mathrm{eff}} = 3.0 \to 4.1$ at $N = 50$--$500$). CDT has a spectral gap $20\times$ larger than causal sets and a stronger entanglement-connectivity correlation ($r = 0.98$ vs.\ $0.85$). However, both discretizations exhibit GUE-like level spacing statistics ($\langle r \rangle \approx 0.60$). We prove that CDT's low $c_{\mathrm{eff}}$ follows from a Kronecker product structure: the antisymmetrized causal matrix decomposes as $C^T - C = A_T \otimes J$, giving $n_{\mathrm{pos}} = \lfloor T/2 \rfloor$ positive eigenvalues (depending only on the number of time slices $T$, not $N$). This dimensional reduction, absent in causal sets, is the mechanism for the central charge discrepancy. The result is fragile: $5\%$ within-slice disorder doubles $c_{\mathrm{eff}}$. This is the first cross-approach comparison of the SJ vacuum and identifies the time foliation as the key structural difference.
\end{abstract}

\maketitle

\section{Introduction}

Discrete approaches to quantum gravity require placing quantum fields on discrete spacetime structures. Two major frameworks---causal set theory~\cite{bombelli1987,surya2019} and causal dynamical triangulations (CDT)~\cite{ambjorn2001,ambjorn2005}---provide complementary discretizations of Lorentzian spacetime. The Sorkin-Johnston (SJ) vacuum~\cite{johnston2009,sorkin2012} defines a canonical quantum state for a free scalar field on any causal set, determined entirely by the causal structure. While the SJ vacuum has been studied extensively on causal sets~\cite{sorkin_yazdi2018,saravani2014}, and field theory on CDT has been explored via lattice methods~\cite{ambjorn2008}, no direct comparison of the SJ vacuum across these two discretizations has been performed.

We fill this gap by computing the SJ vacuum on both 2D causal sets and 2D CDT configurations using identical numerical code. This enables an apples-to-apples comparison of entanglement entropy, spectral gaps, level statistics, and entanglement-geometry correlations, isolating the effects of the discretization scheme from the field theory construction.

\section{Setup}

\subsection{The SJ vacuum}

For a causal set $\mathcal{C}$ with $N$ elements and causal matrix $C_{ij}$ ($= 1$ if $i \prec j$), the Pauli-Jordan function is $i\Delta_{ij} = (2/N)(C_{ji} - C_{ij})$. The SJ Wightman function is the positive-frequency part: $W = \sum_{\lambda_k > 0} \lambda_k |\mathbf{v}_k\rangle\langle\mathbf{v}_k|$, where $H = i \cdot (i\Delta)$ has eigenvalues $\{\lambda_k\}$~\cite{johnston2009}. Entanglement entropy of sub-region $A$ is $S(A) = -\mathrm{Tr}[W_A \ln W_A + (I - W_A)\ln(I - W_A)]$. The same construction applies to any directed acyclic graph with a causal interpretation, including CDT.

\subsection{Causal sets}

We use random 2-orders---intersections of two random total orders on $N$ elements---which faithfully embed into 2D Minkowski spacetime~\cite{glaser2018}. System sizes $N = 20$--$500$.

\subsection{CDT}

We use 2D CDT with cylindrical topology: $T$ time slices, each a ring of spatial edges, connected by $(2,1)$ and $(1,2)$ triangles~\cite{ambjorn2001}. The Regge action is $S_R = \lambda_2 N_2$ (counting triangles). MCMC sampling via Alexander moves at several values of the cosmological constant $\lambda_2$. CDT configurations are converted to a causal matrix by defining $C_{ij} = 1$ iff vertex $j$ is in the causal future of vertex $i$, where the causal structure is determined by the time foliation: $i$ precedes $j$ iff $t_i < t_j$ and there exists a path of triangle edges connecting $i$ to $j$ that moves forward in time at each step. Within a single time slice, vertices are spacelike-separated. System sizes $N = 38$--$127$ (vertices).

\section{Results}

\subsection{Central charge}

The entanglement entropy of a half-system scales as $S(N/2) = (c/3)\ln N + \text{const}$ in 1+1D conformal field theory~\cite{calabrese2004}, with $c = 1$ for a free massless scalar. The effective central charge $c_{\mathrm{eff}} = 3S(N/2)/\ln N$ provides a direct test.

\begin{table}[h]
\centering
\begin{tabular}{lccc}
\toprule
& $N$ & $c_{\mathrm{eff}}$ & Trend \\
\midrule
CDT & 38 & 1.36 & \\
CDT & 83 & 0.98 & $\to 1$ \\
CDT & 127 & 1.15 & \\
\midrule
Causet & 50 & 3.01 & \\
Causet & 100 & 3.32 & \\
Causet & 200 & 3.66 & $\to \infty$ \\
Causet & 300 & 3.86 & \\
Causet & 500 & 4.13 & \\
\bottomrule
\end{tabular}
\caption{Effective central charge $c_{\mathrm{eff}} = 3S(N/2)/\ln N$. CDT converges toward $c = 1$ (free scalar); causal sets diverge.}
\label{tab:central_charge}
\end{table}

CDT yields $c_{\mathrm{eff}} \approx 1.0$--$1.4$, fluctuating around or slightly above $c = 1$ without systematic growth. The three CDT data points do not definitively establish convergence---they fluctuate non-monotonically---but they are clearly in the $c \sim 1$ regime, consistent with one free scalar. Causal sets yield $c_{\mathrm{eff}} \approx 3$ at small $N$, increasing monotonically to $4.1$ at $N = 500$ with no sign of convergence. The causet result persists for MCMC-sampled continuum-phase 2-orders (at $\beta = 0.7\beta_c$), confirming that the divergence is intrinsic to the SJ vacuum on 2-orders, not an artifact of non-manifold configurations. Larger CDT system sizes ($N > 200$) would strengthen the convergence claim.

\subsection{Spectral gap}

The spectral gap of the Pauli-Jordan operator---the smallest positive eigenvalue---measures the ``mass gap'' of the discrete field theory. We normalize by $N$ to obtain the dimensionless gap $\Delta_{\mathrm{gap}} = N \cdot \lambda_{\min}^+$.

\begin{table}[h]
\centering
\begin{tabular}{lcc}
\toprule
Discretization & $\Delta_{\mathrm{gap}} = N \cdot \lambda_{\min}^+$ & Interpretation \\
\midrule
CDT & $56 \pm 8$ & Gapped \\
Causet & $2.8 \pm 0.4$ & Gapless \\
\bottomrule
\end{tabular}
\caption{Dimensionless spectral gap at comparable system sizes ($N \approx 50$). CDT is $20\times$ more gapped than causal sets.}
\label{tab:gap}
\end{table}

CDT has a spectral gap $20\times$ larger than causal sets at comparable system sizes. This suggests that CDT's lattice regularity suppresses the near-zero modes that inflate the causet central charge. The causet gap $\Delta_{\mathrm{gap}} \approx 2.8$ (approximately $N$-independent) is consistent with gapless behavior in the continuum limit, while the CDT gap $\Delta_{\mathrm{gap}} \approx 56$ indicates a lattice-scale mass scale.

\subsection{Level spacing statistics}

Both discretizations exhibit GUE-like level spacing statistics (Table~\ref{tab:lsr}), consistent with quantum chaotic spectral correlations.

\begin{table}[h]
\centering
\begin{tabular}{lcc}
\toprule
Discretization & $\langle r \rangle$ & Class \\
\midrule
CDT & $0.58 \pm 0.03$ & GUE \\
Causet & $0.56 \pm 0.02$ & GUE \\
GUE reference & 0.60 & --- \\
\bottomrule
\end{tabular}
\caption{Level spacing ratio. Both discretizations are GUE-like.}
\label{tab:lsr}
\end{table}

The GUE universality class is shared across discretizations, indicating that quantum chaos is a robust property of the positive-frequency projection on discrete Lorentzian structures, independent of the specific discretization scheme.

\subsection{Entanglement-connectivity correlation}

The ER=EPR-inspired correlation between $|W_{ij}|$ and causal connectivity $\kappa_{ij}$ (number of shared ancestors plus descendants) is stronger on CDT than on causal sets (Table~\ref{tab:epr}).

\begin{table}[h]
\centering
\begin{tabular}{lc}
\toprule
Discretization & $r(|W|, \kappa)$ \\
\midrule
CDT & $0.98 \pm 0.01$ \\
Causet & $0.85 \pm 0.03$ \\
\bottomrule
\end{tabular}
\caption{Entanglement-connectivity correlation. CDT has a nearly perfect linear relationship.}
\label{tab:epr}
\end{table}

CDT's lattice regularity produces a nearly deterministic relationship between $|W|$ and $\kappa$, while the randomness of the causal set sprinkling introduces scatter.

\subsection{Cosmological constant independence}

On CDT, the effective central charge is approximately independent of the cosmological constant coupling $\lambda_2$ over the range tested, confirming that $c_{\mathrm{eff}}$ probes the field theory content (one massless scalar) rather than the gravitational background.

\section{Discussion}

\textbf{1. Different universality classes.} Despite discretizing the same continuum spacetime (2D flat Minkowski), CDT and causal sets yield qualitatively different SJ vacua. CDT recovers $c = 1$; causal sets do not. This means the two discretizations belong to different universality classes for scalar field theory on discrete Lorentzian structures.

\textbf{2. The causet normalization problem.} Our SJ construction uses a $2/N$ normalization for the Pauli-Jordan function, following the convention that ensures $W$ eigenvalues lie in $[0,1]$. Johnston~\cite{johnston2009} uses a different normalization ($1/2$), and Sorkin-Yazdi~\cite{sorkin_yazdi2018} propose truncating near-zero modes. The $2/N$ convention produces too many near-zero modes on causal sets. These modes contribute to the entanglement entropy without corresponding to physical degrees of freedom, inflating $c_{\mathrm{eff}}$. CDT's lattice regularity naturally suppresses these spurious modes. The Sorkin-Yazdi truncation scheme~\cite{sorkin_yazdi2018} may resolve this by removing near-zero modes, but has not been tested at large $N$ on 2-orders.

\textbf{3. Shared quantum chaos.} The GUE level statistics are universal across both discretizations and robust to $N = 500$ on causal sets. This suggests that GUE universality is a property of the positive-frequency projection on \emph{any} discrete causal structure, not specific to either discretization. This is the most robust finding of the cross-approach comparison.

\textbf{4. CDT as benchmark.} CDT's recovery of $c = 1$ validates it as a correct discretization of the 2D free scalar. This makes CDT a useful benchmark: any causal set observable that agrees with CDT is likely physical, while discrepancies (like $c_{\mathrm{eff}}$) point to discretization artifacts.

\textbf{5. Implications for causal set QFT.} The central charge discrepancy suggests that the standard SJ vacuum construction on causal sets, while mathematically well-defined, does not faithfully reproduce the continuum QFT in the simplest case ($c = 1$ free scalar in 2D). Resolving this---whether through modified normalization, mode truncation, or a different vacuum prescription---is an important open problem for causal set field theory.

\textbf{6. The spectral mechanism for $c = 1$ vs.\ $c \to \infty$: a Kronecker product theorem.} The central charge discrepancy has a sharp spectral signature: at $N = 80$, CDT has only 4 positive eigenvalues of $i\Delta$ that are significantly nonzero, while causal sets have 38. We prove that this is a consequence of the time-foliation structure. For a uniform CDT with $T$ time slices of $s$ vertices each ($N = Ts$), the antisymmetrized causal matrix decomposes as a Kronecker product:
\begin{equation}
C^T - C = A_T \otimes J_{s \times s}\,,
\label{eq:kronecker}
\end{equation}
where $A_T$ is the $T \times T$ antisymmetric sign matrix ($[A_T]_{ij} = \mathrm{sign}(j - i)$) and $J_{s \times s}$ is the $s \times s$ all-ones matrix. Since $J$ has rank 1, the eigenvalues of $C^T - C$ are the eigenvalues of $A_T$ scaled by $s$, plus $(s-1)T$ zero eigenvalues. The number of positive eigenvalues of $i\Delta$ is therefore
\begin{equation}
n_{\mathrm{pos}} = \lfloor T/2 \rfloor\,,
\end{equation}
depending \emph{only} on the number of time slices $T$, not on the total number of elements $N$. This was verified with 100\% accuracy across 32 uniform CDT configurations ($T = 3$--$20$, $s = 3$--$10$, $N = 9$--$200$) and 5 non-uniform MCMC-sampled configurations.

Since $T \sim \sqrt{N}$ for typical CDT configurations, $n_{\mathrm{pos}} \sim \sqrt{N}/2$, while causal sets have $n_{\mathrm{pos}} \approx N/2$. The time foliation reduces the $N$-dimensional spectral problem to a $T$-dimensional one, giving $c_{\mathrm{eff}} = O(1)$ on CDT versus $c_{\mathrm{eff}} \to \infty$ on causal sets. The central charge discrepancy is thus not merely a normalization issue but a direct consequence of the dimensional reduction imposed by the time foliation.

\textbf{7. Fragility of CDT's $c = 1$.} While CDT robustly produces $c \approx 1$ for perfectly time-sliced configurations, introducing just 5\% within-slice disorder (randomly reassigning 5\% of vertex time labels within each slice) doubles $c_{\mathrm{eff}}$ from $1.1$ to $2.1$. This indicates that CDT's $c = 1$ is fragile and depends on the strict time-foliation structure, not merely on the causal set being ``manifold-like.'' This fragility may explain why causal sets---which lack a preferred foliation---fail to reproduce $c = 1$.

\textbf{8. Structure vs.\ density: thinning CDT.} To test whether CDT's low $c_{\mathrm{eff}}$ is due to its regular structure or simply its lower density (fewer relations per element), we randomly thin CDT configurations to 30\% of their vertices while preserving the causal structure. The thinned CDT still has $c_{\mathrm{eff}} = 1.0$--$1.3$, confirming that CDT's $c = 1$ comes from its layered \emph{structure}, not its density. This rules out the hypothesis that causal sets' high $c_{\mathrm{eff}}$ is simply an artifact of having more causal relations per element.

\section{Conclusion}

The first cross-approach comparison of the SJ vacuum reveals that CDT reproduces the expected $c = 1$ central charge for a free scalar while causal sets yield a diverging $c_{\mathrm{eff}}$. The mechanism is a Kronecker product decomposition~(\ref{eq:kronecker}): CDT's time foliation reduces the spectral problem from $N$ to $T$ dimensions, giving $n_{\mathrm{pos}} = \lfloor T/2 \rfloor \sim \sqrt{N}$ active modes versus $n_{\mathrm{pos}} \approx N/2$ for causal sets. This result is fragile---$5\%$ within-slice disorder doubles $c_{\mathrm{eff}}$---but robust to thinning (removing $70\%$ of vertices preserves $c \approx 1$). Both discretizations share GUE-like spectral statistics, identifying this as a universal property of discrete Lorentzian field theory. These results establish CDT as a benchmark for causal set field theory and identify the time foliation as the structural feature responsible for correct field theory content.

\begin{acknowledgments}
This work was conducted with assistance from Claude (Anthropic). The SJ vacuum construction follows Johnston~\cite{johnston2009} and Sorkin~\cite{sorkin2012}. The CDT implementation follows Ambj\o{}rn, Jurkiewicz, and Loll~\cite{ambjorn2001}.
\end{acknowledgments}

\begin{thebibliography}{99}

\bibitem{bombelli1987}
L.~Bombelli, J.~Lee, D.~Meyer, and R.~D.~Sorkin,
``Space-time as a causal set,''
Phys.\ Rev.\ Lett.\ \textbf{59}, 521 (1987).

\bibitem{surya2019}
S.~Surya,
``The causal set approach to quantum gravity,''
Living Rev.\ Rel.\ \textbf{22}, 5 (2019),
\href{https://arxiv.org/abs/1903.11544}{arXiv:1903.11544}.

\bibitem{ambjorn2001}
J.~Ambj\o{}rn, J.~Jurkiewicz, and R.~Loll,
``Dynamically triangulating Lorentzian quantum gravity,''
Nucl.\ Phys.\ B \textbf{610}, 347 (2001),
\href{https://arxiv.org/abs/hep-th/0105267}{arXiv:hep-th/0105267}.

\bibitem{ambjorn2005}
J.~Ambj\o{}rn, J.~Jurkiewicz, and R.~Loll,
``Spectral dimension of the universe,''
Phys.\ Rev.\ Lett.\ \textbf{95}, 171301 (2005),
\href{https://arxiv.org/abs/hep-th/0505113}{arXiv:hep-th/0505113}.

\bibitem{johnston2009}
S.~Johnston,
``Feynman propagator for a free scalar field on a causal set,''
Phys.\ Rev.\ Lett.\ \textbf{103}, 180401 (2009),
\href{https://arxiv.org/abs/0909.0944}{arXiv:0909.0944}.

\bibitem{sorkin2012}
R.~D.~Sorkin,
``On the entropy of the vacuum outside a horizon,''
in \emph{General Relativity and Gravitation} (2014), p.~1691,
\href{https://arxiv.org/abs/1205.2953}{arXiv:1205.2953}.

\bibitem{sorkin_yazdi2018}
R.~D.~Sorkin and Y.~K.~Yazdi,
``Entanglement entropy in causal set theory,''
Class.\ Quant.\ Grav.\ \textbf{35}, 074004 (2018),
\href{https://arxiv.org/abs/1611.10281}{arXiv:1611.10281}.

\bibitem{saravani2014}
M.~Saravani, S.~Aslanbeigi, and R.~D.~Sorkin,
``Spacetime entanglement entropy in 1+1 dimensions,''
Class.\ Quant.\ Grav.\ \textbf{31}, 214006 (2014),
\href{https://arxiv.org/abs/1311.7146}{arXiv:1311.7146}.

\bibitem{ambjorn2008}
J.~Ambj\o{}rn, A.~G\"orlich, J.~Jurkiewicz, and R.~Loll,
``Planckian birth of a quantum de Sitter universe,''
Phys.\ Rev.\ Lett.\ \textbf{100}, 091304 (2008),
\href{https://arxiv.org/abs/0712.2485}{arXiv:0712.2485}.

\bibitem{glaser2018}
L.~Glaser, D.~O'Connor, and S.~Surya,
``Finite size scaling in 2d causal set quantum gravity,''
Class.\ Quant.\ Grav.\ \textbf{35}, 024002 (2018),
\href{https://arxiv.org/abs/1706.06432}{arXiv:1706.06432}.

\bibitem{calabrese2004}
P.~Calabrese and J.~Cardy,
``Entanglement entropy and quantum field theory,''
J.\ Stat.\ Mech.\ \textbf{0406}, P06002 (2004),
\href{https://arxiv.org/abs/hep-th/0405152}{arXiv:hep-th/0405152}.

\end{thebibliography}

\end{document}
